\section{Writing a lecture}\label{sec:lectures}

A lecture is a \file{LectureNotes} document. The class comes in two sizes from
one source: the default 14\,pt \emph{handout} (printed or posted for students),
and a 17\,pt \emph{projector} copy for the room, built automatically as a view
(section~\ref{sec:build-and-find}). A typical master, from the demo repo's
\file{derivatives/lectures/1.1-IntroToDerivatives.tex}:

\begin{manexsource}
%! views: projector, instructor
\documentclass{LectureNotes}
\collapseworkspace

% Per-course text-result vocabulary (e.g. APEX's "Key Idea").
\input{../../course-vocab}

\begin{document}
\section*{Introduction to Derivatives}
...
\end{document}
\end{manexsource}

Three preamble notes. The \texttt{\%! views:} comment declares which extra PDFs
the build script produces (here a projector copy and an instructor copy). The
\cs{collapseworkspace} switch selects natural-height boxes -- the right mode when the
notes carry no reserved whitespace; see section~\ref{sec:overlay} for the
whole story. And \file{course-vocab.tex} is \cs{input} \emph{in the preamble}
(its box declarations must run before the document starts), via a relative
path from the repo root -- section~\ref{sec:extending} explains what belongs
in that file.

\subsection{Text results: boxes with declared numbers}

\file{CourseBoxes} provides the boxed results -- and its numbering model is
the one thing to internalise: box numbers are \emph{declared, not counted}.
You type the number, matching the course text, and \cs{label}/\cs{ref} work as
usual:

\begin{manexample}
\begin{definition}{2.1.1}[Derivative]\label{def:deriv}
  The derivative of $f$ at $x$ is
  $f'(x) = \lim_{h \to 0} \frac{f(x+h) - f(x)}{h}$,
  provided the limit exists.
\end{definition}

\begin{theorem}{2.1.4}[Power Rule]\label{thm:power}
  For any real constant $n$, $\lzo{x}\,x^{n} = n\,x^{n-1}$.
\end{theorem}

Theorem~\ref{thm:power} rests on Definition~\ref{def:deriv}; the
text's 2.1.2 and 2.1.3 are simply skipped.
\end{manexample}

Because nothing is counted, the notes mirror the textbook's own numbering and
may skip freely -- the gap above is the feature, not a bug. Keeping LaTeX's
automatic counters aligned with a textbook means fighting them document-wide;
declaring the number is one argument at the call site.
Section~\ref{sec:extending} covers the mechanism and how to add
course-specific vocabulary (\env{keyidea}, \env{axiom}, \dots).

The universal environments, always available: \env{theorem},
\env{proposition}, \env{lemma}, \env{corollary}, \env{definition},
\env{example}, and \env{remark}, all with the same
\texttt{\{number\}[optional name]} signature.

\subsection{Pedagogical asides}

Asides are the \emph{author's} voice -- unnumbered, never referenced, and
deliberately not \cs{label}-able (they are commentary, not results). Two of
the five, live:

\begin{manexample}
\begin{caution}
  The Power Rule is for a \emph{variable base, constant
  exponent}. It does not apply to $2^{x}$.
\end{caution}

\begin{tip}[surviving]
  Remember: $\frac{1}{a+b} \neq \frac{1}{a} + \frac{1}{b}$.
\end{tip}
\end{manexample}

The full set: \env{caution}, \env{takeaway}, \env{yourturn}, \env{prepwork},
and \env{tip} -- where \env{tip} takes an optional audience tier,
\texttt{[surviving]} or \texttt{[thriving]}, that stamps the box for the
struggling or the confident reader respectively (a bare \env{tip} is for
everyone).

\subsection{Importing problems as examples}

The same problem files that grade as ``Q'' on a test
(section~\ref{sec:assessments}) read as worked examples in notes. The class
renders an imported problem with a run-in \textbf{Example.} heading --
unnumbered by default, because the default example is improvised board work,
not a text citation:

\begin{manexsource}
\importproblem{../problems/}{power-rule-basic.tex}
% -> "Example.  Differentiate f(x) = ..."

\importproblem[number={2.1.8}, desc={from APEX}]
  {../problems/}{power-rule-basic.tex}
% -> "Example 2.1.8 (from APEX).  Differentiate f(x) = ..."
\end{manexsource}

The \texttt{number=} key sets a \emph{literal} display number (typically the
textbook's -- the same declared-number philosophy as the boxes; a \cs{ref} to
the problem reads ``2.1.8'' too), and \texttt{desc=} adds a parenthetical
note. Points keys are simply ignored in notes -- the same import line can
move between a test and a lecture unchanged.

Two more import forms round out the lecture repertoire:

\begin{itemize}
  \item \cs{importcheckin} (same arguments as \cs{importproblem}) wraps an
    imported multiple-choice problem in a framed ``CHECK IN:'' box with
    coloured A--E flashcard labels -- the in-class poll format. Same file,
    third appearance.
  \item \env{problemgroup} assembles independent one-file problems under a
    shared stem; in notes the stem stays an unnumbered ``Example.''\ and the
    members letter themselves (a), (b) (see section~\ref{sec:assessments}
    for the graded version).
\end{itemize}

The ``Example'' noun itself is a per-course choice (texts say Example /
Exercise / Problem): redefine \cs{NotesProblemLabelWord} in
\file{course-vocab.tex} and every lecture follows, without touching the fixed
``Q'' on assessments.

\subsection{Instructor-only annotations}

Lectures carry notes-to-self and co-teacher asides that students must never
see. Three shapes, all in slate and all invisible outside an instructor build
(the flag is flipped locally here to show them):

\begin{manexample}
\instructornotestrue
Some lecture prose.
\InstructorInline{Most will try to memorise this before
understanding it -- resist that.}

\begin{InstructorBox}[marking]
  Budget a full board example on expanding-then-cancelling
  before the Power Rule.
\end{InstructorBox}
\end{manexample}

The \cs{InstructorInline} note is a run-in aside in the paragraph flow;
\env{InstructorBox} is a side-barred block (the rotated word says \emph{note},
or \emph{marking} with the optional argument); \cs{InstructorMargin} drops a
bold sans note in the margin (not shown live -- this manual's example boxes
sit too far from the page margin for it to land sensibly). In an instructor
build, every page also gets a page-edge frame reading \emph{Instructor copy}
-- the copy identifies itself even from a glance at a single page.

\subsection{Board space and the projector copy}

\cs{workspace}\texttt{[N]} reserves $N$ lines of handwriting room -- in
lectures, space to improvise at the board while students follow on the
handout. It is the same command assessments use for answer space;
section~\ref{sec:overlay} explains it (and how solutions can \emph{overlay}
that space without moving anything).

The projector view needs no source changes at all: \texttt{\%! views:
projector} makes the build script produce \file{...-PROJECTOR.pdf} alongside
the handout -- 17\,pt type, tighter margins, no page number, and a grey
back-to-top arrow in the footer (the room copy is scrolled, not paged). Two
class behaviours worth knowing: every \cs{subsection} starts a fresh page (a
lecture section is a screenful), and only \cs{section} is numbered.
